\documentclass{article}
\usepackage{booktabs}
\title{FooNet: A Small but Honest Example Paper}
\author{Anti-Autoresearch Eval Fixture}

\begin{document}
\maketitle

\begin{abstract}
We present FooNet, a method for the BarBench task. FooNet reaches
78.0\% accuracy, improving from a 73.1\% baseline to 78.0\% accuracy,
a 6.7\% relative improvement. We report the mean over 5 seeds with
standard deviation, and compare against the strongest published baseline.
\end{abstract}

\section{Introduction}
Prior work established strong baselines on BarBench \cite{smith2024bar}.
We build on the architecture of \cite{lee2023foundations} and add a
lightweight routing module. Our contribution is a single, well-scoped
improvement evaluated on one benchmark.

\section{Method}
FooNet adds a routing module on top of the backbone of
\cite{lee2023foundations}. No test-time labels are used; all calibration
is performed on the training split only.

\section{Experiments}
We evaluate on BarBench. All numbers are the mean of 5 seeds.

\begin{table}[t]
\caption{Main results on BarBench (accuracy, \%, mean of 5 seeds).}
\begin{tabular}{lc}
\toprule
Method & Accuracy \\
\midrule
Baseline \cite{smith2024bar} & 73.1 \\
FooNet (ours) & 78.0 \\
\bottomrule
\end{tabular}
\end{table}

On BarBench, FooNet improves accuracy from 73.1\% to 78.0\%. The gain is
modest and we make no claim of broad generality: results are limited to a
single benchmark with five seeds.

\section{Conclusion}
FooNet yields a 78.0\% accuracy on BarBench. Future work should test
additional benchmarks before any broader claim is warranted.

\end{document}
